\subsection{极大环面与同态的提升}

设 G 是连通紧李群，T 是 G 的极大环面。考虑 G 的换位子群 \(\mathrm{D}(\mathrm{G})\) 及其泛覆盖 \(\tilde{\mathrm{D}}(\mathrm{G})\) ；设 \(p: \tilde{\mathrm{D}}(\mathrm{G}) \to \mathrm{G}\) 是典则态射 \(\tilde{\mathrm{D}}(\mathrm{G}) \to \mathrm{D}(\mathrm{G})\) 与 \(\mathrm{D}(\mathrm{G}) \to \mathrm{G}\) 的复合。则 \(\tilde{\mathrm{D}}(\mathrm{G})\) 是连通紧李群（ \(\S1\) ，第 4 节，命题 4 的推论 2）；此外，T 在 p 下的逆像 \(\tilde{T}\) 是 \(\tilde{\mathrm{D}}(\mathrm{G})\) 的极大环面（第 3 节，命题 1）。

引理 2. 设 H 是李群，\(f_{\mathrm{T}}: \mathrm{T} \to \mathrm{H}\) 与 \(\tilde{f}: \tilde{\mathrm{D}}(\mathrm{G}) \to \mathrm{H}\) 是李群态射，满足 \(f_{\mathrm{T}}(p(t)) = \tilde{f}(t)\) 对一切 \(t \in \tilde{\mathrm{T}}\) 成立。则存在唯一的李群态射 \(f: \mathrm{G} \to \mathrm{H}\) 使得 \(f \circ p = \tilde{f}\) ，且 \(f\) 在 \(\mathrm{T}\) 上的限制为 \(f_{\mathrm{T}}\) 。

令 \(Z = \mathrm{C}(\mathrm{G})_{0}\) ；由 §1 第 4 节命题 4 的推论 1，李群态射 \(g: Z \times \tilde{\mathrm{D}}(\mathrm{G}) \to \mathrm{G}\) （它满足 \(g(z, x) = z^{-1} p(x)\) ）是一个覆盖；它的核由偶 \((z, x)\) （满足 \(p(x) = z\) ）组成，对这些偶有 \(x \in p^{-1}(\mathrm{Z}) \subset \tilde{\mathrm{T}}\) 。由于态射 \((z, x) \mapsto f_{\mathrm{T}}(z^{-1}) \tilde{f}(x)\) （从 \(Z \times \tilde{\mathrm{D}}(\mathrm{G})\) 到 H）把 \(\operatorname{Ker} g\) 映为 \(\{e\}\) ，故存在从 G 到 H 的态射 f 使得 \(f \circ p = \tilde{f}\) 且 \(f(z) = f_{\mathrm{T}}(z)\) 对 \(z \in Z\) 成立。也有 \(f(t) = f_{\mathrm{T}}(t)\) 对 \(t \in p(\tilde{\mathrm{T}})\) 成立；由于 \(T = Z.p(\tilde{\mathrm{T}})\) ，f 在 T 上的限制确为 \(f_{T}\) 。

命题 8. 设 G 是连通紧李群，T 是 G 的极大环面，H 是李群，\(\varphi : \mathrm{L}(\mathrm{G}) \to \mathrm{L}(\mathrm{H})\) 是李代数同态。则存在李群态射 \(f: \mathrm{G} \to \mathrm{H}\) 满足 \(\mathrm{L}(f) = \varphi\) ，当且仅当存在李群态射 \(f_{\mathrm{T}}: \mathrm{T} \to \mathrm{H}\) 满足 \(\mathrm{L}(f_{\mathrm{T}}) = \varphi | \mathrm{L}(\mathrm{T})\) ；此时 \(f_{\mathrm{T}} = f | \mathrm{T}\) 。

若 \(f: \mathrm{G} \to \mathrm{H}\) 是满足 \(\mathrm{L}(f) = \varphi\) 的李群态射，则限制 \(f_{\mathrm{T}}\) （即 \(f\) 在 \(\mathrm{T}\) 上的限制）是从 \(\mathrm{T}\) 到 \(\mathrm{H}\) 的唯一满足 \(\mathrm{L}(f_{\mathrm{T}}) = \varphi | \mathrm{L}(\mathrm{T})\) 的态射。反之，设 \(f_{\mathrm{T}}: \mathrm{T} \to \mathrm{H}\) 是满足 \(\mathrm{L}(f_{\mathrm{T}}) = \varphi | \mathrm{L}(\mathrm{T})\) 的李群态射。设 \(\tilde{\mathrm{D}}(\mathrm{G})\) 与 \(p\) 如上；映射 \(\mathrm{L}(p)\) 诱导一个从 \(\mathrm{L}(\tilde{\mathrm{D}}(\mathrm{G}))\) 到导代数 \(\mathfrak{b}\) （\(\mathrm{L}(\mathrm{G})\) 的）上的同构。存在李群态射 \(\tilde{f}: \tilde{\mathrm{D}}(\mathrm{G}) \to \mathrm{H}\) 使得 \(\mathrm{L}(\tilde{f}) = (\varphi | \mathfrak{b}) \circ \mathrm{L}(p)\) （第 III 章，§6，第 1 节，定理 1）。态射 \(t \mapsto \tilde{f}(t)\) 与 \(t \mapsto f_{\mathrm{T}}(p(t))\) （都是从 \(\tilde{\mathrm{T}}\) 到 \(\mathrm{H}\) ）诱导相同的李代数同态，因而相合。应用引理 2，我们推出存在态射 \(f: \mathrm{G} \to \mathrm{H}\) 使得 \(\mathrm{L}(f)\) 与 \(\varphi\) 在 \(\mathrm{L}(\mathrm{T})\) 与 \(\mathfrak{b}\) 上相合。由于 \(\mathrm{L}(\mathrm{G}) = \mathfrak{b} + \mathrm{L}(\mathrm{T})\) ，我们有 \(\mathrm{L}(f) = \varphi\) 。

命题 9. 设 G 是连通紧李群，T 是 G 的极大环面，H 是李群，\(f: G \to H\) 是态射。则 f 单射当且仅当它在 T 上的限制单射。

事实上，由第 2 节定理 2，G 的正规子群 \(\mathrm{Ker}f\) 约化为单位元当且仅当它与 T 的交约化为单位元。

